\clearpage
\section{Off cone incidence
localization}\label{j-off-cone-incidence-localization}

Technical derivation for review. Equation and subsection numbers in this
appendix are local to the source derivation. Historical local labels are
retained, including numbering gaps or nonsequential labels. The
accompanying source notes retain the original expressions for
comparison.

\subsection{1. A slightly regularized commuting
tree}\label{j-a-slightly-regularized-commuting-tree}

Write a traceless ordered commuting configuration as
\(\lambda =(\lambda _{1},\ldots ,\lambda _{N})\), with
\(\rho ^{2}=\sum  \vert \lambda _{i}\vert ^{2}\). For a proper subset
\(B\) with at least two members put

% DISPLAY j.1
% SOURCE c_B=|B|^(-1) sum_(i in B) lambda_i,
% SOURCE a_B^2=sum_(i in B)|lambda_i-c_B|^2,
% SOURCE d_B^(-2)=sum_(j outside B) (|lambda_j-c_B|^2+a_B^2)^(-1).
\begin{gather*}
\begin{aligned}
& c_{B}=\vert B\vert ^{-1} \sum _{i \in  B} \lambda _{i}, \\
& a_{B}^{2}=\sum _{i \in  B}\vert \lambda _{i}-c_{B}\vert ^{2}, \\
& d_{B}^{-2}=\sum _{j \text{ outside } B} (\vert \lambda _{j}-c_{B}\vert ^{2}+a_{B}^{2})^{-1}.
\end{aligned}\notag
\end{gather*}

The extra \(a_{B}^{2}\) in the denominator is deliberate. It removes
inverse-matrix singularities from the off-cone extension whenever
\(a_{B}>0\), without changing the clustering mechanism. For
\(0<\kappa <1/20, L>1\) and \(\delta =\kappa  \exp (-L)\), use

% DISPLAY j.2
% SOURCE theta_B=Theta(log(a_B/(delta d_B))/L),
% SOURCE C_B=cos(theta_B), S_B=sin(theta_B),
\begin{gather*}
\begin{aligned}
& \theta _{B}=\Theta (\log (a_{B}/(\delta  d_{B}))/L), \\
& C_{B}=\cos (\theta _{B}), S_{B}=\sin (\theta _{B}),
\end{aligned}\notag
\end{gather*}

where \(\Theta\) is a fixed smooth nondecreasing function, \(0\) on
(-infinity,0{]} and \(\pi /2\) on {[}1,infinity). Every derivative of
the trigonometric factors is flat at the constant-tail boundaries. If
\(a_{B}=0\) and all exterior points are separated, \(C_{B}=1\) locally.
Coincident unused nodes are handled by the same identically zero-prefix
mechanism as in the decision tree.

The regularization retains a uniform derivative bound

% DISPLAY j.3
% SOURCE |d a_B|<=1, |d d_B|<=C_N,
% SOURCE |d theta_B|^2 <= C_N ||Theta'||_infinity^2 L^(-2) a_B^(-2)
\begin{gather*}
\begin{aligned}
& \vert d a_{B}\vert \le 1, \vert d d_{B}\vert \le C_{N}, \\
& \vert d \theta _{B}\vert ^{2} \le  C_{N} \Vert \Theta '\Vert _{\infty}^{2} L^{-2} a_{B}^{-2}
\end{aligned}\notag
\end{gather*}

on \(\delta  d_{B}<a_{B}<\kappa  d_{B}\). The constant is independent of
\(\delta\) and \(L\). For example a conservative \(C_{N}\) can be
obtained by taking \(\vert d d_{B}\vert \le 4N\) and then squaring
\(1+4N \kappa\).

For crossing subsets B,C, the same three-point argument gives

% DISPLAY j.4
% SOURCE d_B <= 2a_B+2a_C, d_C <= 2a_B+2a_C.
\begin{gather*}
\begin{aligned}
& d_{B} \le  2a_{B}+2a_{C}, d_{C} \le  2a_{B}+2a_{C}.
\end{aligned}\notag
\end{gather*}

Thus \(C_{B} C_{C}=0\) if \(\kappa <1/4\). All simultaneously selected
clusters are laminar. The decreasing-cardinality decision tree and its
exact pathwise IMS identity therefore still apply.

\subsubsection{Symmetric raw leaf
formula}\label{j-symmetric-raw-leaf-formula}

For an unordered partition \(\pi\) into at least two final blocks, its
tree leaf multiplier is exactly

% DISPLAY j.5
% SOURCE q_pi = product_(B in pi, |B|>=2) C_B
% SOURCE        product_(B proper union of >=2 pi-blocks) S_B.       (1)
\begin{gather*}
\begin{aligned}
& q_{\pi} = \prod _{B \in  \pi , \vert B\vert \ge 2} C_{B} \\
& \prod _{B \text{ proper union of } \ge 2 \pi\text{-blocks}} S_{B}.
\end{aligned}\tag{1}
\end{gather*}

This is not a claim that all visited nodes are unions on their entire
domain. Rather, any omitted rejected crossing node has \(S_{B}=1\) on a
neighborhood of the nonzero final-block selection support, by
laminarity. Nodes strictly inside a selected final block are skipped.
Thus (1) equals the original leaf product as a smooth function,
including its derivatives. In particular

% DISPLAY j.6
% SOURCE sum_pi q_pi^2=1,
% SOURCE sum_pi |d q_pi|^2 = sum_pi q_pi^2 E_pi,                  (2)
\begin{gather*}
\begin{aligned}
& \sum _{\pi} q_{\pi}^{2}=1, \\
& \sum _{\pi} \vert d q_{\pi}\vert ^{2} = \sum _{\pi} q_{\pi}^{2} E_{\pi},
\end{aligned}\tag{2}
\end{gather*}

where \(E_{\pi}\) sums \(\vert d \theta _{B}\vert ^{2}\) for the final
nontrivial blocks and the proper unions in (1). Terms that were skipped
or crossing contribute zero on that leaf. No arbitrary ordering of
equal-cardinality clusters remains in (1).

\subsection{2. Intrinsic matrix extension on one stationary
flag}\label{j-intrinsic-matrix-extension-on-one-stationary-flag}

Use the notation and associated-bundle coordinates of the intrinsic
multiblock geometry note. On an ordered flag of ranks
\(n_{1},\ldots ,n_{k}\), write

% DISPLAY j.7
% SOURCE X=Z+A+Y, B_Z^*Y=0, alpha_a^2=sum_i ||A_i^(a)||_F^2,
% SOURCE r_ab=|z_a-z_b|, r_*=min r_ab.
\begin{gather*}
\begin{aligned}
& X=Z+A+Y, B_{Z}^{*}Y=0, \alpha _{a}^{2}=\sum _{i} \Vert A_{i}^{(a)}\Vert _{F}^{2}, \\
& r_{ab}=\vert z_{a}-z_{b}\vert , r_{*}=\min  r_{ab}.
\end{aligned}\notag
\end{gather*}

For a union I of flag blocks define

% DISPLAY j.8
% SOURCE n_I=sum_(a in I)n_a,
% SOURCE c_I=n_I^(-1)sum_(a in I)n_a z_a,
% SOURCE a_I^2=sum_(a in I)alpha_a^2 + sum_(a in I)n_a|z_a-c_I|^2,
% SOURCE H_(b,I)=sum_i(A_i^(b)+(z_b-c_I)_i I)^2+a_I^2 I,
% SOURCE d_I^(-2)=sum_(b outside I)Tr H_(b,I)^(-1).                (3)
\begin{gather*}
\begin{aligned}
& n_{I}=\sum _{a \in  I}n_{a}, \\
& c_{I}=n_{I}^{-1}\sum _{a \in  I}n_{a} z_{a}, \\
& a_{I}^{2}=\sum _{a \in  I}\alpha _{a}^{2} + \sum _{a \in  I}n_{a}\vert z_{a}-c_{I}\vert ^{2}, \\
& H_{b,I}=\sum _{i}(A_{i}^{(b)}+(z_{b}-c_{I})_{i} I)^{2}+a_{I}^{2} I, \\
& d_{I}^{-2}=\sum _{b \text{ outside } I}\operatorname{Tr}  H_{b,I}^{-1}.
\end{aligned}\tag{3}
\end{gather*}

For a single nontrivial final block, \(a_{I}=\alpha _{a}\). Formula (3)
ignores \(Y\), as it must for the desired leading slow localization. It
uses traces, positive matrix inverses, and projector compressions; it
chooses no internal eigenbasis. If \(a_{I}>0\) every \(H\) is positive.
If \(a_{I}=0\) for a selected final block, exterior separation keeps
every \(H\) positive and its angle has a constant select tail. A union
of at least two blocks has \(a_{I}>0\) when \(r_{*}>0\).

The inverse-trace differentiation estimate remains uniform. If
\(H=\sum _{i} T_{i}^{2}+a^{2}\) I, then Cauchy-Schwarz and
\(\sum _{i}\Vert T_{i} H^{-2}\Vert _{F}^{2}\le \operatorname{Tr}  H^{-3}\)
bound the derivative of \(\operatorname{Tr}  H^{-1}\). The inequality
\(\sum  \operatorname{eigenvalues} (H)^{-3}\le (\sum  \operatorname{eigenvalues} (H)^{-1})^{3}\)
gives a dimension-controlled Lipschitz bound for \(d\). Center
subtraction only introduces the fixed factors from \(n_{I}\) and \(N\).
Hence the \(\theta\) derivative estimate of Section \(1\) holds for the
canonical \((Z,A)\) metric.

For any gauge-invariant scalar f(Z,A), the exact co-metric identity
gives

% DISPLAY j.9
% SOURCE |grad_X f|^2 <= (1+C_N sigma^2)|d_(Z,A)f|^2              (4)
\begin{gather*}
\begin{aligned}
& \vert \nabla _{X} f\vert ^{2} \le  (1+C_{N} \sigma ^{2})\vert d_{Z,A}f\vert ^{2}
\end{aligned}\tag{4}
\end{gather*}

on a thin fast tube. Indeed put \(p=(d_{Z} f,d_{A} f,0)\). Solving the
full coordinate differential and scalar gauge constraint gives the
ambient covector \((p-K^{*} \beta ,\beta )\), where
\(\beta =-M^{-*}O_{H}^{*}p\) and
\(O_{H}^{*}p=O_{C}^{*}d_{Z} f+O_{A}^{*}d_{A} f\). Both \(O_{C}\) and
\(O_{A}\) are projections of \(B_{Y}\). Factoring \(M^{-*}\) with
\(F^{-*}\) on the right gives
\(\vert \beta \vert \le C\vert Y\vert  \vert p\vert /r_{*}\). Since
\(\Vert K\Vert \le 2\vert Y\vert /r_{*}\), the full squared covector
norm is at most \((1+C_{N} \sigma ^{2})\vert p\vert ^{2}\). This
argument includes the full orbital contribution and has no unrelated
internal-radius loss.

\subsection{3. Support margins and an explicit center-gap
bound}\label{j-support-margins-and-an-explicit-center-gap-bound}

If \(q_{\pi}\) is nonzero, all nontrivial final blocks satisfy
\(\alpha _{a}<\kappa  d_{a}\). From

% DISPLAY j.10
% SOURCE d_a^2 <= r_ab^2+alpha_b^2+alpha_a^2
\begin{gather*}
\begin{aligned}
& d_{a}^{2} \le  r_{ab}^{2}+\alpha _{b}^{2}+\alpha _{a}^{2}
\end{aligned}\notag
\end{gather*}

and the analogous inequality with a,b interchanged, obtain

% DISPLAY j.11
% SOURCE alpha_a+alpha_b <= beta r_ab,
% SOURCE beta=2 kappa/sqrt(1-2 kappa^2).                          (5)
\begin{gather*}
\begin{aligned}
& \alpha _{a}+\alpha _{b} \le  \beta  r_{ab}, \\
& \beta =2 \kappa /\sqrt{1-2 \kappa ^{2}}.
\end{aligned}\tag{5}
\end{gather*}

Singleton radii are zero. Choose the tree \(\kappa\) so \(\beta\) is
strictly smaller, for example by a factor four, than the endpoint
constant allowed by the geometric lemma.

Every proper union I of at least two final blocks has been rejected, so
\(a_{I}\ge \delta  d_{I}\) on \(q_{\pi}\)\textquotesingle s support.
Suppose the currently included centers have diameter \(D\). Equation (5)
implies

% DISPLAY j.12
% SOURCE a_I <= sqrt(N(1+beta^2)) D.
\begin{gather*}
\begin{aligned}
& a_{I} \le  \sqrt{N(1+\beta ^{2})} D.
\end{aligned}\notag
\end{gather*}

From \(d_{I}\le a_{I}/\delta\), at least one eigenvalue of some exterior
\(H_{b,I}\) is \(\le N a_{I}^{2}/\delta ^{2}\). Evaluating its unit
eigenvector and using the triangle inequality yields

% DISPLAY j.13
% SOURCE |z_b-c_I| <= sqrt(N) a_I/delta + alpha_b.
\begin{gather*}
\begin{aligned}
& \vert z_{b}-c_{I}\vert  \le  \sqrt{N} a_{I}/\delta  + \alpha _{b}.
\end{aligned}\notag
\end{gather*}

Using \(\alpha _{b}\le \beta \vert z_{b}-z_{a}\vert\) for an included a
and adjoining this center grows \(D\) by at most \(M_{N}/\delta\), where

% DISPLAY j.14
% SOURCE M_N=[N sqrt(1+beta^2)+1]/(1-beta).
\begin{gather*}
\begin{aligned}
& M_{N}=[N \sqrt{1+\beta ^{2}}+1]/(1-\beta ).
\end{aligned}\notag
\end{gather*}

Starting from a closest pair and making at most k-2 additions proves

% DISPLAY j.15
% SOURCE r_* >= Delta_N R,
% SOURCE R^2=r_pi^2+sum_a alpha_a^2,
% SOURCE Delta_N=[N(1+beta^2)]^(-1/2)(delta/M_N)^(N-2).            (6)
\begin{gather*}
\begin{aligned}
& r_{*} \ge  \Delta _{N} R, \\
& R^{2}=r_{\pi}^{2}+\sum _{a} \alpha _{a}^{2}, \\
& \Delta _{N}=[N(1+\beta ^{2})]^{-1/2}(\delta /M_{N})^{N-2}.
\end{aligned}\tag{6}
\end{gather*}

Since \(\rho _{X}^{2}=R^{2}+\vert Y\vert ^{2}\) and
\(\vert Y\vert \le \sigma  r_{*}\), this also gives

% DISPLAY j.16
% SOURCE r_* >= Delta_N rho_X/sqrt(1+2sigma^2).                  (7)
\begin{gather*}
\begin{aligned}
& r_{*} \ge  \Delta _{N} \rho _{X}/\sqrt{1+2\sigma ^{2}}.
\end{aligned}\tag{7}
\end{gather*}

All nonzero raw weights, and the closures of their derivative supports,
have this margin. This is the key quantitative substitute for an
unspecified finite cover.

For a selected-block angle transition,

% DISPLAY j.17
% SOURCE d_a >= (1-beta) r_*/sqrt(N),
% SOURCE alpha_a >= delta(1-beta)r_*/sqrt(N).                    (8)
\begin{gather*}
\begin{aligned}
& d_{a} \ge  (1-\beta ) r_{*}/\sqrt{N}, \\
& \alpha _{a} \ge  \delta (1-\beta )r_{*}/\sqrt{N}.
\end{aligned}\tag{8}
\end{gather*}

For a proper-union angle, \(a_{I}\) dominates its center-relative
radius. Thus all nonconstant cutoff factors live at a fixed positive
multiple of \(\rho\), once \(N,\kappa ,L\) have been fixed. Internal
coincidences never cause an unbounded derivative of a surviving
multiplier.

\subsection{4. All branches, finite count, and smooth zero
extension}\label{j-all-branches-finite-count-and-smooth-zero-extension}

For every ordered composition \(N=n_{1}+\ldots +n_{k}, k\ge 2\), take
the incidence set of stationary ordered orthogonal projector flags for
\(\Phi _{X}\). Admit branches inside a padded domain with:

\begin{itemize}
\tightlist
\item
  endpoint bound strictly larger than \(\beta\), but still inside the
  stated geometric regime
\item
  center gap strictly smaller than the lower bound (7), for example one
  fourth of that bound
\item
  fast width strictly larger than the support of an invariant smooth
  fast cutoff
\end{itemize}

Use a smooth cutoff \(\chi _{Y}=1\) for \(\vert Y\vert \le \sigma\) h/4
and \(\chi _{Y}=0\) for \(\vert Y\vert \ge \sigma\) h/2, with
\(h=(\sum _{a<b}r_{ab}^{-2})^{-1/2}\). The geometric domain can allow
\(\vert Y\vert <\sigma  r_{*}\). Put

% DISPLAY j.18
% SOURCE b_j=(k!)^(-1/2) chi_Y q_pi                             (9)
\begin{gather*}
\begin{aligned}
& b_{j}=(k!)^{-1/2} \chi _{Y} q_{\pi}
\end{aligned}\tag{9}
\end{gather*}

on each stationary flag sheet \(j\). Extend it by zero outside the
padded support. Equations (5)-(7), the flat cutoff tails, and the strict
fast margin make this a smooth finite-cover section. There is no cutoff
associated with a choice of coordinate chart or block frame.

The Hessian estimate from the geometry note makes each admitted
stationary flag a nondegenerate critical point. In particular it is
locally a smooth branch, and no branch can be born or disappear inside
the nonzero support. Branches can have monodromy; chart transitions
merely permute direct-sum components.

\subsubsection{Explicit N-only multiplicity
bound}\label{j-explicit-n-only-multiplicity-bound}

A crude bound is sufficient. A flag of type \((n_{a})\) has real
dimension \(m=N^{2}-\sum  n_{a}^{2}\le N^{2}\). It is covered by at most
\(N\)! pivot charts. In a complex block-lower-unipotent chart, every
cumulative orthogonal projector is

% DISPLAY j.19
% SOURCE V(V^*V)^(-1)V^*.
\begin{gather*}
\begin{aligned}
& V(V^{*}V)^{-1}V^{*}.
\end{aligned}\notag
\end{gather*}

A common denominator for all projectors has degree at most \(2N^{2}\) in
real chart coordinates. \(\Phi\) has numerator and denominator degree at
most \(4N^{2}\); clearing denominators from each first derivative gives
polynomial equations of degree at most \(8N^{2}\). Nondegenerate real
critical points are isolated complex zeros as well. The elementary
Bezout isolated-zero bound therefore gives at most \((8N^{2})^{m}\) per
chart. Summing the at most \(2^{N-1}\) ordered compositions gives the
conservative bound

% DISPLAY j.20
% SOURCE D_N=2^(N-1) N! (8N^2)^(N^2).                           (10)
\begin{gather*}
\begin{aligned}
& D_{N}=2^{N-1} N! (8N^{2})^{N^{2}}.
\end{aligned}\tag{10}
\end{gather*}

The bound is deliberately huge and only used to ensure constants do not
secretly depend on the logarithmic inner scale. It counts ordered
sheets; (9) corrects their actual multiplicity.

\subsubsection{Why commuting admitted sheets are exactly cluster
partitions}\label{j-why-commuting-admitted-sheets-are-exactly-cluster-partitions}

Let \(B(X)=\sum _{i<j}\Vert [X_{i},X_{j}]\Vert _{F}^{2}\). On a
stationary thin flag, the transverse linear commutator has norm

% DISPLAY j.21
% SOURCE ||L_Z Y||^2=sum_(a<b) r_ab^2 |Y_ab|^2
\begin{gather*}
\begin{aligned}
& \Vert L_{Z} Y\Vert ^{2}=\sum _{a<b} r_{ab}^{2} \vert Y_{ab}\vert ^{2}
\end{aligned}\notag
\end{gather*}

with the corresponding fixed real-block normalization. Endpoint
perturbations are bounded by \(C\) times the padded geometric endpoint
constant times this norm (and by \(C \beta\) on the raw-weight support),
and the quadratic \(Y\) term by \(C \sigma\) times it. With the fixed
small geometric constants,

% DISPLAY j.22
% SOURCE B(X)^(1/2) >= c_N (sum r_ab^2|Y_ab|^2)^(1/2).          (11)
\begin{gather*}
\begin{aligned}
& B(X)^{1/2} \ge  c_{N} (\sum  r_{ab}^{2}\vert Y_{ab}\vert ^{2})^{1/2}.
\end{aligned}\tag{11}
\end{gather*}

Consequently a commuting admitted sheet has \(Y=0\). Its projectors
group common eigenspaces; a repeated joint eigenvalue cannot be split
between two selected thin blocks, by (5). Each unordered particle
partition appears exactly \(k\)! times among all ordered compositions
and flags. Thus on the commuting cone

% DISPLAY j.23
% SOURCE Q=sum_j b_j^2=1.                                      (12)
\begin{gather*}
\begin{aligned}
& Q=\sum _{j} b_{j}^{2}=1.
\end{aligned}\tag{12}
\end{gather*}

No extra \(N\)! multiplier belongs in (9).

Equation (11) also ensures that on a sufficiently small near-commuting
region the fast cutoff is identically one on every contributing sheet.
Using (7) and \(h\ge r_{*}/\sqrt{k(k-1)/2}\), it suffices to take

% DISPLAY j.24
% SOURCE t=B(X)/rho_X^4 <= c_N sigma^2 Delta_N^4.                (13)
\begin{gather*}
\begin{aligned}
& t=B(X)/\rho _{X}^{4} \le  c_{N} \sigma ^{2} \Delta _{N}^{4}.
\end{aligned}\tag{13}
\end{gather*}

The constant in (13) can be reduced to include every fixed geometric
margin. There is then no fast-cutoff derivative in the near-cone
localization; outside it, its usual cost could instead be paid by a
retained fast reserve.

\subsection{5. The exact cone identity and the small off-cone
defect}\label{j-the-exact-cone-identity-and-the-small-off-cone-defect}

Where \(Q>0\) define \(w_{j}=b_{j}/\sqrt{Q}\). Then
\(\sum _{j} w_{j}^{2}=1\). Let

% DISPLAY j.25
% SOURCE E(X)=sum_j |grad_X w_j|^2,
% SOURCE R(X)=sum_j w_j^2 E_j(X),                               (14)
\begin{gather*}
\begin{aligned}
& E(X)=\sum _{j} \vert \nabla _{X} w_{j}\vert ^{2}, \\
& R(X)=\sum _{j} w_{j}^{2} E_{j}(X),
\end{aligned}\tag{14}
\end{gather*}

where \(E_{j}\) is the sum of the ambient
\(\vert \nabla  \theta _{I}\vert ^{2}\) for the single-block select
factors and proper-union reject factors of that sheet. All expressions
are taken before any slow/fast projection.

At a commuting \(X\) on an admitted sheet, \(Y=0\). A pointwise common
diagonalization verifies that the first derivative of every scalar in
(3) in an off-diagonal internal direction is zero. This is a calculation
of an invariant Frechet derivative, not a choice of smoothly varying
eigenvectors. It remains true at repeated joint eigenvalues: trace
derivatives of inverse matrices against off-diagonal perturbations
vanish, and the radius-zero select factor is constant on a neighborhood.
Further, \(K=O_{C}=O_{A}=0\) at \(Y=0\) in the exact kinetic identity.
Therefore the ambient gradients of all raw weights equal their
commuting-base gradients, including at singular strata. The same
statement holds for the angles wherever their derivatives matter.

The whole symmetric commuting identity (2) therefore implies

% DISPLAY j.26
% SOURCE Q=1, dQ=0, E=R                                      (15)
\begin{gather*}
\begin{aligned}
& Q=1, dQ=0, E=R
\end{aligned}\tag{15}
\end{gather*}

at every nonzero commuting tuple. This is stronger than just equality of
the weight values.

On the unit sphere, all supported branch sums and their first two
derivatives are bounded: the incidence Hessian is uniformly invertible
after (7), the flag multiplicity is bounded by (10), all nonconstant
cutoff transition scales \(a_{I}\) obey fixed positive lower bounds, and
the support stays strictly inside the padded branch domain. Accordingly
\(Q\) and E-R are smooth scalar functions in a neighborhood of the
compact commuting unit locus. There is a finite explicitly defined
derivative bound \(K_{N,\kappa ,L,\sigma ,\Theta}\), obtained as the
supremum of their first derivatives on a closed smaller padded
neighborhood, such that

% DISPLAY j.27
% SOURCE |Q(X)-1|+|E(X)-R(X)| <= K dist(X,K_comm)                (16)
\begin{gather*}
\begin{aligned}
& \vert Q(X)-1\vert +\vert E(X)-R(X)\vert  \le  K \operatorname{dist} (X,K_{\mathrm{comm}})
\end{aligned}\tag{16}
\end{gather*}

for unit \(X\) close enough to that locus. This constant can grow with
\(\Delta _{N}^{-1}\); it is not asserted uniform in the shape gap. Its
role is to determine a later near-cone thickness, never to appear as an
unpaid inverse-square coefficient.

For every \(\epsilon >0\) choose a neighborhood thickness \(\eta >0\)
small enough that \(Q\ge 1/2\) and \(K \eta \le \epsilon\). Homogeneity
then proves

% DISPLAY j.28
% SOURCE E(X) <= R(X)+epsilon rho_X^(-2)                        (17)
\begin{gather*}
\begin{aligned}
& E(X) \le  R(X)+\epsilon  \rho _{X}^{-2}
\end{aligned}\tag{17}
\end{gather*}

on that conical neighborhood. If desired, the elementary estimate of
Section \(6\) makes the choice in terms of \(t\): take
\(t_{1}\le (\eta /(2 C_{N}))^{4}\), together with (13) and the
branch-continuation radius. This is a concrete order of choices, not a
simultaneous assumption \(\kappa =\kappa (\Delta _{N})\).

\subsubsection{Why direct leafwise normalization was
insufficient}\label{j-why-direct-leafwise-normalization-was-insufficient}

For a single select factor,
\(\vert d \cos  \theta \vert ^{2}=\sin ^{2} \theta  \vert d \theta \vert ^{2}\)
cannot be bounded by \(\cos ^{2} \theta\) times the same cost near its
zero. Hence one cannot bound \(\vert dw_{j}\vert ^{2}\) by
\(w_{j}^{2} E_{j}\) separately. Nor may one charge every nested
transition to every final leaf: a tiny subcluster transition inside a
larger selected block is not controlled with a shape-independent
coefficient by the latter\textquotesingle s total-radius Hardy estimate.
Only the summed defect argument (15)-(17) repairs this issue.

\subsection{6. An elementary almost-commuting distance
certificate}\label{j-an-elementary-almost-commuting-distance-certificate}

There is no need to assume a smooth retraction onto the singular
commuting cone. For a traceless Hermitian nine-tuple \(X\) of norm one,
put
\(\epsilon _{0}=\max _{i<j}\Vert [X_{i},X_{j}]\Vert _{F}\le \sqrt{t}\).
An elementary spectral-block construction produces a commuting traceless
tuple \(Z\) with

% DISPLAY j.29
% SOURCE ||X-Z|| <= C_N sqrt(epsilon_0) <= C_N t^(1/4),
% SOURCE C_N=729 sqrt(N^3+9)                                   (18)
\begin{gather*}
\begin{aligned}
& \Vert X-Z\Vert  \le  C_{N} \sqrt{\epsilon _{0}} \le  C_{N} t^{1/4}, \\
& C_{N}=729 \sqrt{N^{3}+9}
\end{aligned}\tag{18}
\end{gather*}

as a conservative choice.

To see this, process one component at a time inside the already
constructed invariant blocks. In each current block, group consecutive
eigenvalues across gaps at most \(\eta =\sqrt{\epsilon _{j}}\). Distinct
groups have spectral distance greater than \(\eta\). Replace the
processed component on each new spectral block by its trace-average; its
tuple error is at most \(N^{3/2} \eta\). Delete interblock entries of
all remaining components. The Frobenius norm of each deleted part is at
most \(\epsilon _{j}/\eta\), by its commutator with the processed
component. The new block-diagonal commutator differs from the compressed
old commutator only by two products of deleted off-block parts. Hence
\(\epsilon _{j+1}\le \epsilon _{j}+2(\epsilon _{j}/\eta )^{2}\le 3\epsilon _{j}\).
At most nine steps give (18); all changes preserve traces. Previously
processed components are scalar on their blocks and stay commuting with
subsequent steps.

This proof uses spectral subspaces only to establish a distance
inequality. The actual localization never uses this approximate tuple,
its ordering, or any internal eigenbasis. If the right side of (18) is
at most \(1/2\), normalize \(Z\) to norm one; its distance from \(X\) is
at most \(2C_{N} t^{1/4}\). Thus (16) can be converted into an explicit
\(t^{1/4}\) modulus with the finite derivative constant \(K\).

\subsection{7. Payment and physical
measure}\label{j-payment-and-physical-measure}

The first term \(R\) in (17) has the original leafwise payments:

\begin{enumerate}
\def\labelenumi{\arabic{enumi}.}
\tightlist
\item
  A selected-block transition costs at most
  \(C_{N} L^{-2} \alpha _{a}^{-2}\), with
  \(\alpha _{a}\ge c_{N} \delta  \Delta _{N} r_{\pi}\) by (8). On the
  physical low coefficient the lower-rank soft Hardy estimate pays this
  from an arbitrarily small internal-form fraction, with
  faster-than-inverse-square exterior defect.
\item
  A proper-union transition has
  \(a_{I}^{2}=\sum  \alpha _{a}^{2}+s_{I}^{2}\), where \(s_{I}\) is the
  canonical relative-center radius of that union. Therefore
  \(a_{I}^{-2}\le s_{I}^{-2}\), paid by ordinary Hilbert-valued
  relative-center Hardy in dimension 9(l-1). There are at most \(2^{N}\)
  such terms per sheet. Their coefficient \(C_{N}/L^{2}\) can be made
  arbitrarily small before \(\delta\) and \(\Delta _{N}\) are fixed.
\item
  On the colored complementary coefficient all scalar costs are bounded
  by \(C_{N,\kappa ,L} r_{\pi}^{-2}\), using (6)-(8), and are paid by a
  retained fast gap at sufficiently large radius. The physical
  lower-rank Hardy theorem is not applied to colored vectors.
\item
  The additional \(\epsilon  \rho _{X}^{-2}\) in (17) is at most
  \(\epsilon  r_{\pi}^{-2}\). On the low coefficient it uses an
  arbitrarily small fraction of the center Hardy reserve; on the
  complement it is absorbed by the fast reserve at large radius. If the
  near-cone region is set by \(t\), the existing outer cutoff has its
  separately paid commutator-potential IMS cost.
\end{enumerate}

The only form inputs in this payment are the previously identified
internal soft-Hardy estimates and the separately supplied all-vector
multiblock comparison with sufficient retained reserves. This
construction does not manufacture those reserves.

On each flag sheet the physical density is the exact \(J=\det  M\) from
the intrinsic geometry note. Same-sheet frame changes are unitary
residual-block/normal-frame transitions and introduce no scalar chart
partition. Distinct sheets remain distinct summands. Pull back the
physical section to the incidence cover and multiply by \(w_{j}\). Then

% DISPLAY j.30
% SOURCE ||u||^2=sum_j ||w_j u||^2,
% SOURCE h[u]=sum_j h[w_j u]-integral E(X)|u|^2                  (19)
\begin{gather*}
\begin{aligned}
& \Vert u\Vert ^{2}=\sum _{j} \Vert w_{j} u\Vert ^{2}, \\
& h[u]=\sum _{j} h[w_{j} u]-\int  E(X)\vert u\vert ^{2}
\end{aligned}\tag{19}
\end{gather*}

with sheet integration against the pulled-back physical measure,
equivalently \(J\) dZ dA dY after gauge integration. The \(k\)! factor
in (9) and normalization, rather than discarding branches, enforce
single counting. Local form estimates can therefore be integrated
chart-free on every sheet. No globally smooth projector, angular
gauge-fixing multiplier, or internal bound-state projector is used.

\subsection{8. Noncircular parameter
order}\label{j-noncircular-parameter-order}

For each fixed \(N\) and target loss \(\epsilon\):

\begin{enumerate}
\def\labelenumi{\arabic{enumi}.}
\tightlist
\item
  Fix the necessary fractions of center/internal/fast reserves in the
  conditional local form theorem.
\item
  Choose the geometric endpoint and fast constants; choose tree
  \(\kappa\) so \(\beta\) in (5) lies strictly inside the endpoint
  bound. Any analytic requirement on these constants must itself be
  shape-independent, as required by the local comparison.
\item
  Choose \(L\) so the sum of at most \(2^{N}\) logarithmic costs is paid
  by the allocated slow reserves.
\item
  Set \(\delta =\kappa  e^{-L}, \Delta _{N}\) from (6), padded incidence
  domains, and the finite derivative/\allowbreak{}continuation constants \(K\).
\item
  Choose the near-commuting threshold \(t_{1}\) using (13), (16)-(18),
  so the added IMS defect is at most \(\epsilon  \rho ^{-2}\).
\item
  Choose the outer cutoff parameters and finally the exterior radius
  large enough to pay all fixed-shape constants, complementary
  inverse-square costs, and soft-Hardy remainders.
\end{enumerate}

No rank-uniform statement or effective small radius is asserted. The
point is that every large delta-dependent constant is either multiplied
by a later freely chosen cone thickness or absorbed at a later exterior
radius.
